% ══════════════════════════════════════════════════════════════════
%  FICHE D'EXERCICES — Chapitre 12 : Dosages acido-basiques et solutions tampons
% ══════════════════════════════════════════════════════════════════
\section{Exercices type Baccalauréat}
\textit{Données à \SI{25}{\celsius} : $pK_e = 14$. À l'équivalence d'un dosage
acido-basique : $n_{\text{acide}} = n_{\text{base}}$ (pour un monoacide/monobase).}

\rubrique{Dosages acide fort / base forte}

\exo{}\diff{1} On dose \SI{20}{\milli\litre} d'acide chlorhydrique par une solution
d'hydroxyde de sodium de concentration $C_b = \SI{0.10}{\mol\per\litre}$. L'équivalence
est obtenue pour $V_{bE} = \SI{15}{\milli\litre}$.
\begin{enumerate}[label=\arabic*)]
\item Écrire l'équation de la réaction de dosage.
\item Déterminer la concentration de l'acide.
\item Quel est le pH à l'équivalence ? Justifier.
\end{enumerate}

\exo{}\diff{2} Pour le dosage précédent, indiquer parmi les indicateurs colorés
suivants celui (ceux) qui convien(nen)t : hélianthine (zone $3{,}1$–$4{,}4$),
B.B.T. (zone $6{,}0$–$7{,}6$), phénolphtaléine (zone $8{,}2$–$10$). Justifier.

\rubrique{Dosages acide faible / base forte}

\exo{}\diff{2} On dose \SI{20}{\milli\litre} d'acide éthanoïque
($C_a = \SI{0.10}{\mol\per\litre}$, $pK_a = 4{,}8$) par la soude
$C_b = \SI{0.10}{\mol\per\litre}$.
\begin{enumerate}[label=\arabic*)]
\item Écrire l'équation de la réaction de dosage.
\item Calculer le volume de soude à l'équivalence.
\item À la demi-équivalence, que vaut le pH ? Justifier.
\item Le pH à l'équivalence est-il acide, neutre ou basique ? Justifier et en déduire
l'indicateur coloré approprié.
\end{enumerate}

\exo{}\diff{2} La courbe de dosage d'un acide $AH$ par la soude présente un saut de
pH autour de $V_{bE} = \SI{12}{\milli\litre}$, et à la demi-équivalence le pH vaut
$4{,}2$. La prise d'essai d'acide est de \SI{20}{\milli\litre} et
$C_b = \SI{0.10}{\mol\per\litre}$.
\begin{enumerate}[label=\arabic*)]
\item Déterminer la concentration de l'acide $AH$.
\item Que représente le pH à la demi-équivalence ? En déduire le $pK_a$ du couple.
\item L'acide est-il fort ou faible ? Justifier.
\end{enumerate}

\rubrique{Solutions tampons}

\exo{}\diff{1} Définir une solution tampon. Citer deux de ses propriétés.

\exo{}\diff{2} On veut préparer une solution tampon de pH $= 4{,}8$ à partir d'acide
éthanoïque et d'éthanoate de sodium ($pK_a = 4{,}8$).
\begin{enumerate}[label=\arabic*)]
\item Quelle relation doit exister entre $[\ce{CH3COOH}]$ et $[\ce{CH3COO-}]$ ?
\item On mélange \SI{50}{\milli\litre} d'acide éthanoïque
($\SI{0.10}{\mol\per\litre}$) et \SI{50}{\milli\litre} d'éthanoate de sodium
($\SI{0.10}{\mol\per\litre}$). Calculer le pH du mélange.
\end{enumerate}

\exo{}\diff{2} À la demi-équivalence du dosage d'un acide faible par une base forte,
on obtient une solution tampon. Expliquer pourquoi, et donner la valeur du pH en
fonction du $pK_a$.

\rubrique{Problème de synthèse — type Baccalauréat}

\sujetbac[d'après Baccalauréat D, Cameroun]
On dose un volume $V_a = \SI{10}{\milli\litre}$ d'une solution d'acide éthanoïque par
une solution d'hydroxyde de sodium $C_b = \SI{0.10}{\mol\per\litre}$. On relève :
\begin{center}\small
\begin{tabular}{l|cccccccc}
$V_b$ (\si{\milli\litre}) & 0 & 2 & 5 & 8 & 9{,}5 & 10 & 10{,}5 & 12 \\ \hline
pH & 2{,}9 & 4{,}1 & 4{,}8 & 5{,}4 & 6{,}3 & 8{,}4 & 10{,}6 & 11{,}3 \\
\end{tabular}
\end{center}
\begin{enumerate}[label=\arabic*)]
\item Tracer la courbe $pH = f(V_b)$ et déterminer les coordonnées du point
d'équivalence (méthode des tangentes).
\item En déduire la concentration de l'acide éthanoïque.
\item Déterminer le $pK_a$ du couple à partir de la courbe (demi-équivalence).
\item Justifier le caractère basique de la solution à l'équivalence et choisir un
indicateur coloré adapté.
\item Quelle est la particularité de la solution obtenue à la demi-équivalence ?
\end{enumerate}
